\section{Galaxies, Black Holes, and Our Universe} % Overview
Galaxies, the vast collections of stars, gas, and dust that populate our Universe, exhibit a remarkable diversity in their properties.  One key characteristic used to classify galaxies is their star formation activity, with galaxies broadly categorised as either active star-forming galaxies (SFG) or quiescent galaxies (QG; \citealt{whitaker_newfirm_2011}). This distinction is reflected in their observed colours: SFGs, rich in young, hot, and massive stars, typically appear bluer, while QGs, dominated by older stellar populations, are characterised by redder colours \citep{baldry_quantifying_2004}. The colours of galaxies provide crucial insights into their underlying physics and evolutionary history.  Since a galaxy's light is the integrated emission of its constituent stars and dust, the observed colours reflect the stellar populations and star formation history. Young, massive stars emit predominantly in the spectrum's ultraviolet (UV) portion, giving SFGs their characteristic blue hue \citep{madau_cosmic_2014}.  As star formation declines and these stars evolve or die, the UV emission diminishes, leading to the redder colours observed in QGs. However, interpreting galaxy colours can be complex, with factors such as dust obscuration and cosmological redshift also influencing the observed colours. Dust can absorb and re-emit starlight, shifting the emission towards longer wavelengths, particularly in the mid-infrared (MIR) and far-infrared (FIR; \citealt{madau_cosmic_2014}).  Furthermore, the expansion of the universe causes the light from distant galaxies to be stretched, leading to a redshift that can make intrinsically blue galaxies appear redder.
Adding another layer of complexity is the presence of Active Galactic Nuclei (AGN) – compact, luminous regions at the centres of some galaxies powered by accreting supermassive black holes. The intense radiation from an AGN can significantly contribute to the total light from a galaxy, potentially outshining the host galaxy and altering its observed colours. In particular, this contamination can potentially affect colour-colour diagrams, a diagnostic tool used to compare the brightness of a source at different wavelengths.  Understanding the impact of this `AGN contamination' on galaxy colours and colour-colour diagrams is crucial for accurately interpreting galaxy properties and their evolution. Despite its importance, the precise relationship between AGN and their host galaxies remains an active research area.
\section{Motivation}
This research is driven by two primary motivations: the desire to further our understanding of galaxy evolution and the complex interplay between galaxies and their central supermassive black holes (SMBHs). \\\\
Firstly, we seek to contribute to the ongoing investigation of the relationship between SMBHs and their host galaxies.  Observational evidence suggests a strong connection between the mass of the SMBH and various properties of the host galaxy, such as the velocity dispersion of stars and the luminosity of the central bulge \citep{kormendy_coevolution_2013}.  This suggests a co-evolutionary relationship, where the growth and activity of the SMBH are intimately linked to the evolution of the host galaxy. However, the precise nature of this relationship and the underlying physical mechanisms remain poorly understood.  By investigating the impact of AGN contamination on galaxy photometry, we aim to gain a clearer picture of how AGN activity can influence the observed properties of galaxies and contribute to their evolution.\\\\ Secondly, we aim to improve galaxy colour diagnostics for large photometric surveys by developing methods to quantify the effects of AGN contamination on galaxy colours. This is particularly crucial for future surveys, such as the Evolutionary Map of the Universe Survey (EMU; \citealt{norris_evolutionary_2021}), which will provide a wealth of data that can be used to study galaxy populations and their evolution. However, the accurate interpretation of this data relies on robust diagnostic tools that can effectively classify galaxies and determine their properties. As the introduction highlights, AGN contamination can significantly affect galaxy colours, potentially leading to misclassification and inaccurate property estimations using these diagnostics \citep{cowley_decoupled_2018}. By developing methods to quantify the effects of AGN contamination, we aim to provide a foundation for enabling more robust analyses of large-scale photometric surveys and ultimately leading to a deeper understanding of galaxy evolution.
\section{Thesis Outline}
In addition to this introduction, this thesis comprises eight sections. Section 2 outlines the central research question, defines the study's aims, and presents a testable hypothesis. Section 3 provides a comprehensive review of relevant literature, exploring the relationship between AGN and their host galaxies, including the complexities of spectral energy distribution (SED) modelling and fitting, the use of colour-colour diagnostic tools, and the challenge of AGN contamination. Section 4 details the ZFOURGE survey, the source of the observational data used in this thesis, highlighting its strengths and ancillary datasets. Section 5 outlines the methodology employed, encompassing the theoretical modelling of AGN contamination in various colour spaces, introducing the SED templates and AGN models used, and discussing the observational analysis techniques. Section 6 presents the results of both the theoretical modelling and the observational analysis, including findings from SED fitting and decomposition. In Section 7, we discuss the implications of these results, exploring the effects of AGN on colour-colour diagnostics and considering the broader implications and limitations of this research. Finally, Section 8 summarises the main conclusions and offers recommendations for future research.\\\\
This research operates within a $\Lambda CDM$ cosmological framework, assuming the following parameters: $H_0 = 70 km \ s^{-1} \ Mpc^{-1}$, $\Omega_M = 0.3$, and $\Omega_{\Lambda} = 0.7$.
